\chapter{Implementazione}
\label{chap:implementation}

\section{Il wrapper di Zenoh pico}
\label{sec:zenoh-pico-wrapper}

A differenza della versione standard di Zenoh, che offre una libreria nativa Rust,
Zenoh pico dispone solo di una libreria C, così da poter sfruttare PlatformIO come
strumento per supportare facilmente tutte le famiglie di dispositivi embedded più
popolari.
Per utilizzare questa versione di Zenoh, di conseguenza, è necessario
creare dei \emph{binding} che permettano di interagire con la libreria tramite
chiamate a funzioni inter-linguaggio.

Inoltre, la creazione di un wrapper Rust permette di nascondere le \ac{API} C dietro
delle strutture più adeguate e coerenti con lo stile di programmazione e gestione
della memoria unici del linguaggio.

\subsection{Integrazione della libreria}
\label{ssec:zenoh-pico-library-compilation}

Per semplificare l'integrazione della libreria e per ridurre i tempi di sviluppo,
il contributo di questa tesi si limita i vincoli di compatibilità nel contesto dei
sistemi embedded solo ai dispositivi della famiglia Espressif e compatibili (\ref{constraint:esp-compatibility}).

Questa limitazione permette di sfruttare il framework \ac{ESP-IDF}, introdotto nella
\cref{ssec:esp-idf}, che offre un sistema di componenti
per integrare librerie esterne al processo di compilazione: semplicemente definendo
i file sorgenti da considerare, questi vengono compilati con le impostazioni necessarie
per avere accesso alle funzioni della libreria standard C offerte dalla piattaforma,
per poi essere esposti pubblicamente come modulo all'interno del framework.
Nonostante Zenoh pico sia una libreria pensata per sistemi embedded, la sua compatibilità
con essi è limitata a quelli che dispongono di un'implementazione, anche se limitata,
della libreria standard C; rendendo quindi necessario sfruttare questa funzionalità.

Per poi rendere accessibili le funzioni e tipi del framework originale C,
\ac{ESP-IDF} utilizza \emph{bindgen} \footnote{\url{https://rust-lang.github.io/rust-bindgen/}},
uno strumento per la generazione automatica di bindings Rust per chiamate a funzioni
di librerie esterne.
In questo modo, il programma può utilizzare un'unica libreria per interagire con
il sistema tramite astrazioni di alto livello, utilizzare le primitive C per un controllo
più diretto ed accedere ai componenti aggiunti esposti tramite moduli dedicati (\Cref{lst:esp-idf-ffi}).

\lstinputlisting[
    language=Rust,
    caption={[Esempio di utilizzo del framework \ac{ESP-IDF} da Rust]
    Esempio di utilizzo del framework \ac{ESP-IDF} da Rust: astrazioni di alto
    livello, primitive C e utilizzo di componenti extra tramite i moduli generati
    automaticamente da bindgen.},
    label={lst:esp-idf-ffi}
]{listings/esp-idf-ffi.rs}

\subsection{Le macro in Rust}
\label{ssec:macros-in-rust}

Rust fornisce costrutti di \emph{metaprogrammazione}~\cite{enwiki:1335302027} per
la generazione di codice a tempo di compilazione. Alla base di questo sistema si
trovano le \emph{macro}, un costrutto concettualmente simile a una funzione, che
riceve in input e produce come output dei flussi di simboli (\emph{token stream})
elaborati dal compilatore (\Cref{lst:vec-macro}).

\lstinputlisting[
    language=Rust,
    caption={[Implementazione semplificata della macro \texttt{vec}]
    Implementazione semplificata della macro \texttt{vec}~\cite[Listing 20-35]{rust-book},
    che inizializza un vettore con le espressioni fornite in input.},
    label={lst:vec-macro}
]{listings/vec-macro.rs}

Le macro si dividono in due categorie principali, dichiarative e procedurali:
mentre le prime permettono di definire il codice da generare in maniera
statica rispetto all'input, fornendo solamente qualche costrutto di
\emph{pattern matching} sulla struttura delle espressioni in input,
le macro procedurali sono effettive funzioni che operano su flussi di simboli,
prendendo in input e fornendo come output delle \texttt{TokenStream}
interpretabili dal compilatore \cite[20.5.3]{rust-book}.
Quest'ultime vengono utilizzate come attributi da poter applicare a definizioni di
simboli per modificarne la struttura o il comportamento. Un esempio sono le \emph{derive macro},
un caso particolare di macro procedurali per implementare automaticamente un \emph{trait}
per un tipo di dato.

Siccome la struttura \texttt{TokenStream} è molto primitiva e priva di operazioni
di alto livello per la manipolazione e la generazione dei simboli, nel tempo sono
nate diverse librerie per semplificarne l'utilizzo, come \emph{syn} \footnote{\url{https://docs.rs/syn/2.0.118/syn/}}
e \emph{darling} \footnote{\url{https://docs.rs/darling/0.23.0/darling/}} per il
parsing automatico dell'input in maniera strutturata o \emph{quote} \footnote{\url{https://docs.rs/quote/1.0.45/quote/}}
per la generazione dell'output a partire da codice standard (\Cref{lst:derive-macro-example}).

\lstinputlisting[
    language=Rust,
    caption={Esempio di implementazione di una \emph{derive macro} utilizzando le
    librerie syn, quote e darling.},
    label={lst:derive-macro-example}
]{listings/derive-macro-example.rs}

\subsection{L'ownership model in Zenoh pico}
\label{ssec:ownership-model-in-zenoh-pico}

Come zenoh pico simuli l'ownership model di Rust in C per mantenere un API il più
simile possibile alla libreria originale e per facilitare la gestione della memoria.

\subsection{L'implementazione tramite macro}
\label{ssec:implementation-via-macros}

Come il wrapper sia stato implementato tramite macro sfruttando la strutturazione
simil-rust della libreria originale \cref{ssec:ownership-model-in-zenoh-pico}

\subsection{Il sistema di closure}
\label{sec:closure-system}

Descrizione di come è stato gestito il concetto di closure di zenoh pico e di come
sono stati utilizzati i signal (non del sistema operativo, quelli di embassy) per
permettere un utilizzo asincrono invece che a callback.

\section{Implementazione della rete}
\label{sec:network-implementation}

\subsection{Il trait \texttt{Network}}
\label{sec:the-network-trait}

Spiegazione più dettagliata del concetto di network e dei due metodi da cui è composta.

\section{La ZenohNetwork}
\label{sec:the-zenoh-network}
\sidenote{La generalizzazione per zenoh e zenoh pico devo ancora implementarla, quindi
potrebbe cambiare questa sezione}

Spiegazione del funzionamento della rete zenoh e della generalizzazione riguardo
all'implementazione usata.

\subsection{Il \texttt{AtomicMessagesStore}}
\label{sec:the-atomic-message-store}

Spiegazione del concetto di message store, della gestione della concorrenza e della
lifespan dei nodi in base all'ultimo update.

\subsubsection{La task di hearthbit}
\label{sec:the-heartbeat-task}

Spiegazione della hearthbit task per mantenere vivo un nodo.

\subsection{Implementazione tramite Zenoh pico}
\label{sec:zenoh-pico-implementation}

Approfondimento sull'implementazione zenoh pico

\subsection{Implementazione tramite Zenoh}
\label{sec:zenoh-implementation}

Approfondimento sull'implementazione zenoh
